\documentclass[tikz,border=6pt]{standalone}
\usepackage[no-math]{fontspec}
\usepackage[brazil]{babel}
\usepackage{amsmath,amssymb}
\IfFileExists{C:/Windows/Fonts/Lato_Regular.ttf}{%
  \setmainfont{Lato}[
    Path = C:/Windows/Fonts/,
    UprightFont = Lato_Regular.ttf,
    BoldFont = Lato_Bold.ttf,
    ItalicFont = Lato_Italic.ttf,
    BoldItalicFont = Lato_BoldItalic.ttf
  ]
  \setsansfont{Lato}[
    Path = C:/Windows/Fonts/,
    UprightFont = Lato_Regular.ttf,
    BoldFont = Lato_Bold.ttf,
    ItalicFont = Lato_Italic.ttf,
    BoldItalicFont = Lato_BoldItalic.ttf
  ]
}{%
  \RequirePackage[default]{lato}
}
\renewcommand{\familydefault}{\sfdefault}

\usepackage{xcolor}
\usepackage{tikz}
\usepackage{pgfplots}
\pgfplotsset{compat=1.18}
\usetikzlibrary{arrows.meta,positioning,shapes.geometric,fit,backgrounds,calc,decorations.pathreplacing,patterns}

% Paleta institucional UEFS / CHF662
\definecolor{uefsNavy}{HTML}{1B2A4A}
\definecolor{uefsTeal}{HTML}{0E7C7B}
\definecolor{uefsForest}{HTML}{2D6A4F}
\definecolor{uefsGold}{HTML}{C4983B}
\definecolor{uefsGoldDark}{HTML}{8F6B1E}
\definecolor{uefsAlert}{HTML}{B02A37}
\definecolor{uefsOrange}{HTML}{D95F02}
\definecolor{uefsInk}{HTML}{22252A}
\definecolor{uefsMuted}{HTML}{5C6470}
\definecolor{uefsWater}{HTML}{2E86AB}
\definecolor{uefsSand}{HTML}{D9B382}
\definecolor{bgNavy}{HTML}{F2F5F9}
\definecolor{bgTeal}{HTML}{E9F5F4}
\definecolor{bgForest}{HTML}{EEF5F0}
\definecolor{bgGold}{HTML}{FCF7EB}
\definecolor{bgAlert}{HTML}{FDF2F2}
\definecolor{bgGray}{HTML}{F8FAFC}
\definecolor{lineGray}{HTML}{CFD8DC}

\tikzset{
  >=Stealth,
  aaCard/.style={draw=#1, fill=white, rounded corners=4pt, line width=1pt, inner sep=6pt, align=center, font=\small},
  aaBox/.style={draw=#1, rounded corners=4pt, line width=1.1pt, inner sep=6pt, align=center, font=\footnotesize\sffamily, text=uefsInk},
  aaTitle/.style={font=\normalsize\bfseries\sffamily, text=uefsNavy, align=center},
  aaSub/.style={font=\footnotesize\sffamily, text=uefsMuted, align=center},
  aaArrow/.style={->, >=Stealth, line width=1.4pt},
  aaDArrow/.style={<->, >=Stealth, line width=1.4pt}
}


\begin{document}
\begin{tikzpicture}

% Título
\node[aaTitle] at (6.5, 6.95)
  {Bacia Hidrográfica como Unidade Funcional e Hierarquia Fluvial de Strahler};

% Divisor topográfico suave (bacia hidrográfica)
\fill[bgTeal!75] plot[smooth cycle, tension=0.75] coordinates {
  (0.9, 3.5) (1.7, 5.2) (4.2, 6.0) (7.4, 5.9) (9.8, 4.9)
  (10.8, 3.2) (9.8, 1.7) (6.9, 1.0) (3.4, 1.1) (1.4, 2.0)
};
\draw[dashed, color=uefsGold, line width=2.1pt] plot[smooth cycle, tension=0.75] coordinates {
  (0.9, 3.5) (1.7, 5.2) (4.2, 6.0) (7.4, 5.9) (9.8, 4.9)
  (10.8, 3.2) (9.8, 1.7) (6.9, 1.0) (3.4, 1.1) (1.4, 2.0)
};

% Faixa ripária / APP ao longo do canal principal (3ª ordem)
\draw[uefsForest, opacity=0.24, line width=12pt, line cap=round, line join=round]
  (5.5, 3.6) -- (6.5, 3.4) -- (8.3, 3.28) -- (10.55, 3.21);

% Tributários de 1ª ordem
\draw[uefsWater!70, line width=1.4pt, line cap=round]
  (2.1, 5.0) -- (3.4, 4.35)
  (2.4, 3.7) -- (3.4, 4.35)
  (4.1, 5.6) -- (5.0, 4.65)
  (5.9, 5.55) -- (5.0, 4.65)
  (2.2, 2.2) -- (3.8, 2.6)
  (2.9, 1.4) -- (3.8, 2.6)
  (5.1, 1.95) -- (5.45, 3.08)
  (7.0, 5.1) -- (7.85, 4.15)
  (8.7, 4.9) -- (7.85, 4.15);

% Tributários de 2ª ordem
\draw[uefsWater, line width=2.3pt, line cap=round]
  (3.4, 4.35) -- (5.5, 3.6)
  (5.0, 4.65) -- (5.5, 3.6)
  (3.8, 2.6) -- (6.5, 3.4)
  (7.85, 4.15) -- (8.3, 3.28);

% Canal principal de 3ª ordem
\draw[uefsNavy, line width=3.6pt, line cap=round, line join=round]
  (5.5, 3.6) -- (6.5, 3.4) -- (8.3, 3.28) -- (10.8, 3.2);

% Rótulos de ordem 1 (em zonas abertas fora dos segmentos)
\node[font=\scriptsize\bfseries\sffamily, text=uefsWater] at (1.92, 5.12) {1};
\node[font=\scriptsize\bfseries\sffamily, text=uefsWater] at (2.20, 3.62) {1};
\node[font=\scriptsize\bfseries\sffamily, text=uefsWater] at (3.95, 5.75) {1};
\node[font=\scriptsize\bfseries\sffamily, text=uefsWater] at (6.05, 5.72) {1};
\node[font=\scriptsize\bfseries\sffamily, text=uefsWater] at (2.00, 2.25) {1};
\node[font=\scriptsize\bfseries\sffamily, text=uefsWater] at (2.75, 1.32) {1};
\node[font=\scriptsize\bfseries\sffamily, text=uefsWater] at (4.98, 1.82) {1};
\node[font=\scriptsize\bfseries\sffamily, text=uefsWater] at (6.88, 5.25) {1};
\node[font=\scriptsize\bfseries\sffamily, text=uefsWater] at (8.85, 5.05) {1};

% Rótulos de ordem 2 (fora dos segmentos)
\node[font=\footnotesize\bfseries\sffamily, text=uefsNavy] at (4.45, 4.18) {2};
\node[font=\footnotesize\bfseries\sffamily, text=uefsNavy] at (5.45, 4.22) {2};
\node[font=\footnotesize\bfseries\sffamily, text=uefsNavy] at (4.65, 3.08) {2};
\node[font=\footnotesize\bfseries\sffamily, text=uefsNavy] at (8.30, 3.85) {2};

% Rótulo do canal principal (3ª ordem) em zona aberta abaixo do trecho inicial
\node[font=\footnotesize\bfseries\sffamily, text=uefsNavy, align=center]
  at (7.15, 2.55) {3\textsuperscript{a} ordem (Canal principal)};

% Exutório
\fill[uefsAlert] (10.8, 3.2) circle (4.5pt);
\node[font=\footnotesize\bfseries\sffamily, text=uefsAlert, align=left, anchor=west] (exu)
  at (11.2, 4.45) {Exutório\\[-1pt]{\scriptsize (seção de saída de}\\{\scriptsize água e sedimentos)}};
\draw[aaArrow, color=uefsAlert, line width=1.3pt] (11.8, 3.95) -- (10.93, 3.32);

% Divisor Topográfico
\node[font=\footnotesize\bfseries\sffamily, text=uefsGoldDark, align=left, anchor=west] (div)
  at (8.55, 6.22) {Divisor Topográfico\\[-1pt]{\scriptsize (separa bacias vizinhas)}};
\draw[aaArrow, color=uefsGoldDark, line width=1.3pt] (8.45, 6.18) -- (6.3, 5.98);

% Área Ripária / APP
\node[font=\footnotesize\bfseries\sffamily, text=uefsForest, align=left, anchor=west] (app)
  at (8.35, 1.25) {Área Ripária / APP\\[-1pt]{\scriptsize (filtro de sedimentos)}};
\draw[aaArrow, color=uefsForest, line width=1.3pt] (9.1, 1.65) -- (9.1, 3.05);

% Quadro inferior da Regra de Strahler (sem ponto de exclamação)
\node[aaBox=uefsNavy, fill=white, inner sep=5pt, font=\footnotesize\bfseries\sffamily, text=uefsNavy]
  at (5.8, 0.35)
  {Regra de Strahler:\quad $1 + 1 \rightarrow 2 \;\mid\; 2 + 2 \rightarrow 3 \;\mid\; 2 + 1 \rightarrow 2$ (mantém a maior ordem)};

\end{tikzpicture}
\end{document}
